% ============================================================
% FRONTMATTER — Pagina titolo и copyright
% ============================================================

% --- Pagina titolo ---
\begin{titlepage}
\centering
\vspace*{2cm}

% Linea decorativa superiore
{\color{manning-red}\rule{\textwidth}{2pt}}
\vspace{1.5cm}

{\fontsize{28}{34}\selectfont\sffamily\bfseries\color{manning-dark} Sviluppo 0-Code\par}
\vspace{0.8cm}
{\Large\sffamily\color{manning-gray} Diventa un Architetto del Contesto: \\ 
il metodo ADLC per progettare e rilasciare software reale\\ guidando l'AI\par}
\vspace{2.5cm}
{\large\sffamily Riccardo Rolla\par}
\vfill

% Linea decorativa inferiore
{\color{manning-red}\rule{\textwidth}{2pt}}
\vspace{0.5cm}
{\normalsize\sffamily\color{manning-gray} Prima Edizione — 2026\par}
\end{titlepage}

% --- Pagina copyright ---
\thispagestyle{empty}
\vspace*{\fill}
\noindent\textsf{\textbf{Sviluppo 0-Code}}\\
\textsf{\textit{Diventa un Architetto del Contesto: il metodo ADLC per progettare e rilasciare software reale guidando l'AI}\\[1em]
\noindent Copyright \textcopyright\ 2026 [Riccardo Rolla]\\
Tutti i diritti riservati.\\[1em]
\noindent Prima edizione: Aprile 2026\\[1em]
\noindent I nomi di prodotti, marchi registrati e aziende menzionati in questo libro
sono proprietà dei rispettivi titolari. L'utilizzo in questo testo è esclusivamente
a scopo identificativo e non implica alcuna affiliazione o approvazione.\\[1em]
\noindent Compilato con \LuaLaTeX.
\newpage

% --- Come leggere questo libro ---
\chapter*{Come leggere questo libro}
\addcontentsline{toc}{chapter}{Come leggere questo libro}

Questo libro utilizza diversi riquadri speciali per evidenziare informazioni
importanti durante la lettura. Impara a riconoscerli prima di iniziare:
ogni tipo ha un significato preciso.

\begin{tipbox}
I \textbf{Suggerimenti} offrono consigli pratici, scorciatoie e best practice
che rendono il lavoro più efficiente. Sono frutto dell'esperienza diretta
con gli strumenti descritti nel libro.
\end{tipbox}

\begin{warnbox}
I riquadri \textbf{Attenzione} segnalano insidie comuni, errori frequenti
o comportamenti non intuitivi degli strumenti. Leggerli con cura
ti eviterà problemi.
\end{warnbox}

\begin{notebox}
Le \textbf{Note} forniscono approfondimenti, contesto aggiuntivo o
riferimenti utili. Non sono indispensabili per seguire il testo,
ma arricchiscono la comprensione.
\end{notebox}

\begin{dangerbox}
I riquadri \textbf{Pericolo} avvertono di rischi seri: operazioni
che possono causare perdita di dati, costi imprevisti o problemi
di sicurezza. Non ignorarli mai.
\end{dangerbox}

\begin{versionbox}
Le \textbf{Note di versione} indicano funzionalità legate a versioni
specifiche degli strumenti. L'ecosistema dell'IA evolve rapidamente:
queste note ti aiutano a capire cosa potrebbe essere cambiato.
\end{versionbox}

\begin{costbox}
I riquadri \textbf{Quanto costa?} indicano quando un'operazione potrebbe
comportare costi (API, servizi cloud, licenze). Così puoi decidere
consapevolmente se procedere.
\end{costbox}

\begin{checkpointbox}
I \textbf{Checkpoint} sono momenti di verifica: ti confermano che il tuo
progetto è nello stato corretto prima di procedere al passo successivo.
\end{checkpointbox}

\begin{praticabox}
I riquadri \textbf{Pratica} propongono esercizi guidati per mettere le mani
in pasta. Sono progettati per consolidare i concetti appena presentati.
\end{praticabox}

\medskip
\noindent I \textbf{blocchi di codice} sono evidenziati con sfondo grigio e una
barra laterale. Puoi digitarli manualmente per esercizio, ma ricorda:
il punto del metodo 0-code è che spesso sarà l'IA a generarli per te.

\newpage
